\section{Exterior apolarity and the four-form minimum}
\label{sec:four-form-problem}

Let $K$ be a field, let $V$ be an $m$-dimensional $K$-vector space, and
let $0\ne\omega\in\bigwedge^d V^*$.  For $0\leq j\leq d$, write
\[
  C_j(\omega):\bigwedge^jV\longrightarrow
  \left(\bigwedge^{d-j}V\right)^*,
  \qquad
  C_j(\omega)(a)(b)=\omega(a\wedge b).
\]
The convention is intrinsic and does not require an inner product.  When
$d=4$, $C_2(\omega)$ is symmetric because degree-two exterior elements
commute.  If $\operatorname{char}K\ne2$, it is, up to a nonzero scalar,
the polar form of the Pluecker quadric associated with $\omega$.

Define the homogeneous exterior annihilator and apolar quotient by
\[
  \operatorname{Ann}_\wedge(\omega)_j
  =\ker C_j(\omega),
  \qquad
  A_\omega=\bigwedge V/\operatorname{Ann}_\wedge(\omega),
\]
where every component of degree greater than $d$ is included in the
annihilator.

\begin{proposition}[Exterior apolar duality]
\label{prop:exterior-apolar-duality}
The space $\operatorname{Ann}_\wedge(\omega)$ is a homogeneous ideal.
Multiplication followed by evaluation against $\omega$ induces perfect
pairings
\[
  (A_\omega)_j\times(A_\omega)_{d-j}\longrightarrow K
  \quad(0\leq j\leq d).
\]
Consequently $A_\omega$ is a graded skew-commutative Poincare-duality
algebra of formal dimension $d$, and
\[
  \dim_K(A_\omega)_j=\operatorname{rank}_K C_j(\omega).
\]
In particular, for a concise four-form,
\[
  \operatorname{Hilb}(A_\omega)
  =\bigl(1,m,\operatorname{rank}C_2(\omega),m,1\bigr).
\]
\end{proposition}

\begin{proof}
If $a\in\ker C_j(\omega)$ and $c\in\bigwedge^qV$, then for every
$b\in\bigwedge^{d-j-q}V$ one has
$\omega((a\wedge c)\wedge b)=0$, after only the exterior sign needed to
move $c$ past $b$.  Thus the annihilator is an ideal.  The left kernel of
the degree-$(j,d-j)$ pairing is exactly $\ker C_j(\omega)$ and its right
kernel is exactly $\ker C_{d-j}(\omega)$.  After quotienting, the pairing
is nondegenerate on both sides.  Equivalently, the induced injection from
$(A_\omega)_j$ to $(A_\omega)_{d-j}^*$ is an isomorphism, since $C_j$ and
its signed transpose $C_{d-j}$ have the same rank.  The dimension formula
and the four-form specialization follow.
\end{proof}

This proposition identifies the current extremal quantity as
\[
  \mu_4^K(m)=min\left\{
    \dim_K(A_\omega)_2:
    \omega\in\bigwedge^4(K^m)^*,\ 
    \dim_K(A_\omega)_1=m
  \right\}.
\]
It is important not to replace $A_\omega$ by an arbitrary
Poincare-duality algebra.  In formal dimension three, the multiplication
of a degree-one-generated Poincare-duality algebra is recovered from its
alternating top form~\cite{Suciu2020PoincareDuality}.  In formal dimension
four, an arbitrary algebra may
contain additional degree-two multiplication data; $A_\omega$ is the
canonical exterior apolar quotient attached to the specified four-form.

\subsection{The exact seven-dimensional value}

After a basis is fixed, coordinate Hodge duality identifies covariant
three-forms with covariant four-forms.  It sends general-linear orbits to
general-linear orbits projectively: the acting matrix is changed by the
surjective contragredient automorphism $g\mapsto g^{-T}$ and by a determinant
scalar.  Cohen and Helminck classify the nonzero seven-dimensional
three-form orbits over an algebraically closed field into the nine classes
$f_1,\ldots,f_9$ listed in their Table~1
\cite{CohenHelminck1988TrilinearForms}.  In characteristic zero, coordinate
duality and exact elimination give the following contraction data.

\begin{center}
\begin{tabular}{c|c|c|c}
trivector orbit & representative & $\operatorname{rank}C_1(*f_i)$
& $\operatorname{rank}C_2(*f_i)$ \\
\hline
$f_1$ & $123$ & 4 & 6 \\
$f_2$ & $123+145$ & 6 & 10 \\
$f_3$ & $123+456$ & 7 & 12 \\
$f_4$ & $162+243+135$ & 7 & 12 \\
$f_5$ & $123+456+147$ & 7 & 16 \\
$f_6$ & $152+174+163+243$ & 7 & 15 \\
$f_7$ & $146+157+245+367$ & 7 & 18 \\
$f_8$ & $123+145+167$ & 6 & 15 \\
$f_9$ & $123+456+147+257+367$ & 7 & 21
\end{tabular}
\end{center}

Here $ijk$ abbreviates $e_i^*\wedge e_j^*\wedge e_k^*$; reordering a
displayed term contributes its exterior sign.

\begin{theorem}[Seven-dimensional minimum]
\label{thm:mu-four-seven}
Over $K=\mathbb Q$ and over every algebraically closed field of
characteristic zero,
\[
  \mu_4^K(7)=12.
\]
\end{theorem}

\begin{proof}
Over an algebraically closed characteristic-zero field, the cited
classification is exhaustive.  The table shows that precisely
$f_3,f_4,f_5,f_6,f_7,f_9$ yield concise dual four-forms and that their
smallest middle contraction rank is $12$.  The representative $f_3$ is
rational; its coordinate dual is
\[
  *f_3=e_{4567}^*-e_{1237}^*,
\]
which has contraction-rank vector $(1,7,12,7,1)$ and supplies the rational
upper bound.  Conversely, scalar extension from $\mathbb Q$ to
$\overline{\mathbb Q}$ preserves the ranks of all contraction matrices.
The algebraically closed lower bound therefore also applies over
$\mathbb Q$.
\end{proof}

The JSON artifact
\texttt{seven\_dimensional\_orbit\_ranks.json} records the nine source
representatives and the rank table.  Its independent standard-library
verifier checks the payload hash, exterior signs, coordinate duals, all
five contraction ranks, and the stated minimum.  The verifier does not
claim to re-prove the literature classification; the source theorem is
the separate orbit-coverage input.
